\documentclass[12pt]{article}
%\usepackage{ctex}
%\usepackage[no-math]{fontspec}
\usepackage[top=0.5in,bottom=0.5in,left=0.5in,right=0.5in]{geometry}
\usepackage{xeCJK}[2011/05/20 v2.4.1]
\setCJKmainfont{KaiTi}

\defaultfontfeatures{Mapping=tex-text}

\usepackage{tikz}
\usetikzlibrary{matrix,arrows,decorations.pathmorphing,positioning}
\usepackage{amsmath}
%\usepackage{amssymb}
%\usepackage{graphicx}
%\usepackage{amsfonts}
%\usepackage{listings}
%\usepackage{natbib}
%\usepackage{dcolumn}
%\usepackage{tabularx}
\usepackage{ulem}
%\usepackage{enumitem}
%\usepackage{array}
%\usepackage{dashrule}
\usepackage{scalerel}
%\usepackage[usestackEOL]{stackengine}
%\usepackage{adjustbox}
%\usepackage{multirow}
%\usepackage{booktabs}
%\usepackage[math]{anttor}
%\usepackage{emerald}
%\usepackage[sfdefault]{quattrocento}

% for mathop
\newcommand{\myunit}{0.8cm}
\newcommand{\num}[1]{\node[rectangle,outer sep=1,inner sep=1,minimum width=\myunit,minimum height=\myunit,anchor=south] {#1};}
\newcommand{\nul}{\node[draw,rectangle, rounded corners,outer sep=1,inner sep=1,minimum width=\myunit,minimum height=\myunit]{};}
\newcommand{\nulls}{\node[rectangle, rounded corners,outer sep=1,inner sep=1,minimum width=\myunit,minimum height=\myunit]{};}
\newcommand{\lin}{\node[rectangle,outer sep=1,inner sep=0,minimum width=\myunit*1.1,anchor=north]{};}
\newcommand{\arcs}{\node[fill=black,rectangle,outer sep=0,inner sep=0,minimum height=\myunit*1.1,anchor=south]{};}
%----------------

% \gridraiseamount is a font-specific value
\newcommand\gridraiseamount{0.12em}

% \grid for a single character
\newcommand\grid[1]{%
\begin{tikzpicture}[baseline=(char.base)]
  \path[use as bounding box]
    (0,0) rectangle (1em,1em);
  \draw[help lines,step=0.5em]
    (0,0) grid (1em,1em);
  \draw[help lines,dashed]
    (0,0) -- (1em,1em)  (0,1em) -- (1em,0);
  \node[inner sep=0pt,anchor=base west]
    (char) at (0em,\gridraiseamount) {#1};
\end{tikzpicture}}

% \grid for a single character
\newcommand\engrid[1]{%
\begin{tikzpicture}[baseline=(char.base)]
  \path[use as bounding box]
    (0,0) rectangle (5.4em,1em);
  \draw[help lines,xstep=.6em,ystep=0.5em]
    (0,0) grid (3.6em,1em);
  \draw[help lines,xstep=.4em,ystep=1em]
    (3.6em,0em) grid (4em,1em);
  \draw[help lines,xstep=.8em,ystep=0.5em]
    (4em,0em) grid (4.8em,1em);
  \draw[help lines,xstep=.6em,ystep=0.5em]
    (4.8em,0em) grid (5.4em,1em);
  \node[inner sep=0pt,anchor=base west]
    (char) at (0em,\gridraiseamount) {#1};
\end{tikzpicture}}

% define height of lowercase letters
\newcommand{\lowercaseheight}{1ex}
% define height of uppercase letters
\newcommand{\uppercaseheight}{1.6ex}
% define depth of descenders
\newcommand{\descenderdepth}{-0.5ex}
\newcommand\enlines[1]{\fontsize{42pt}{\baselineskip}\selectfont
\raggedright
\begin{tikzpicture}[overlay]
  % baseline
  \draw ((0,0) -- ++(\textwidth,0);
  % lowercase height
  \draw [ultra thin] (0,\lowercaseheight) -- ++(\textwidth,0);
  % uppercase height
  \draw (0,\uppercaseheight) -- ++(\textwidth,0);
  % descender depth
  \draw [ultra thin] (0,\descenderdepth) -- ++(\textwidth,0);
  \node[inner sep=0pt,anchor=base west]
    (char) at (0em,0em) {#1};
\end{tikzpicture}\\[1.55cm]
}

\pagestyle{empty}
\begin{document}
%{\Large
%\begin{center}
%  \hspace{-.2cm}姓名:{\fontsize{38pt}{\baselineskip}\selectfont\grid{}\grid{}\grid{}}\quad Date:\uline{\quad}/\uline{\quad}/\uline{\quad\quad}\quad Temp:\uline{\quad\quad}$^{\circ}{\mathrm F}$\quad 温度:\uline{\quad\quad}$^{\circ}{\mathrm C}$
%\end{center}}

{\Large
\begin{center}
  \hspace{-.2cm}中文家庭作业(\today)\quad 姓名:{\fontsize{38pt}{\baselineskip}\selectfont\grid{}\grid{}\grid{}}\quad
  %日期:\uline{\quad}/\uline{\quad}/\uline{\quad}
  分数：\uline{\qquad}
\end{center}}

{\fontsize{18pt}{\baselineskip}\selectfont
\xeCJKsetup{PunctStyle=plain}
\vspace{0.5cm}
抄写古诗《赋得古原草送别》两遍}
%填空
{\fontsize{40pt}{\baselineskip}\selectfont
\xeCJKsetup{PunctStyle=plain}
\vspace{0.5cm}
\center
%\grid{} \grid{}\\[.2cm]
\grid{} \grid{} \grid{} \grid{} \grid{} \grid{}\\[.1cm]
\grid{} \grid{} \grid{} \grid{} \grid{} \grid{}\\[.1cm]
\grid{} \grid{} \grid{} \grid{} \grid{} \grid{}\\[.1cm]
\grid{} \grid{} \grid{} \grid{} \grid{} \grid{}\\[.5cm]
%\grid{} \grid{}\\[.2cm]
\grid{} \grid{} \grid{} \grid{} \grid{} \grid{}\\[.1cm]
\grid{} \grid{} \grid{} \grid{} \grid{} \grid{}\\[.1cm]
\grid{} \grid{} \grid{} \grid{} \grid{} \grid{}\\[.1cm]
\grid{} \grid{} \grid{} \grid{} \grid{} \engrid{}\\[.5cm]
}
\clearpage
{\fontsize{18pt}{\baselineskip}\selectfont
\xeCJKsetup{PunctStyle=plain}
\vspace{0.5cm}
写出有相同偏旁的字\\[.8cm]}
%填空
{\fontsize{40pt}{\baselineskip}\selectfont
\xeCJKsetup{PunctStyle=plain}
\vspace{0.5cm}
\grid{行} \grid{} \grid{} \grid{} \grid{} \grid{}\\[.1cm]
\grid{杜} \grid{} \grid{} \grid{} \grid{} \grid{}\\[.1cm]
\grid{牧} \grid{} \grid{} \grid{} \grid{} \grid{}\\[.1cm]
\grid{远} \grid{} \grid{} \grid{} \grid{} \grid{}\\[.1cm]
\grid{寒} \grid{} \grid{} \grid{} \grid{} \grid{}\\[.1cm]
\grid{深} \grid{} \grid{} \grid{} \grid{} \grid{}\\[.1cm]
\grid{停} \grid{} \grid{} \grid{} \grid{} \grid{}\\[.1cm]
\grid{晚} \grid{} \grid{} \grid{} \grid{} \grid{}\\[.1cm]
\grid{霜} \grid{} \grid{} \grid{} \grid{} \grid{}\\[.1cm]
\grid{叶} \grid{} \grid{} \grid{} \grid{} \grid{}\\[.1cm]
\grid{红} \grid{} \grid{} \grid{} \grid{} \grid{}\\[.1cm]
\grid{花} \grid{} \grid{} \grid{} \grid{} \grid{}\\[.1cm]
}
\end{document}
